% ===========================================================
%  APPENDICES
% ===========================================================

\chapter{MJCF File of the Musculoskeletal Model}
\label{app:mjcf}

An excerpt from the MJCF file describing the nine-muscle arm. The muscle force
scales follow \parencite{Holzbaur2005}; the inertial parameters follow
\parencite{Winter2009}. \Cref{tab:muscles} lists the values as compiled. The
complete file is in the Git repository accompanying this thesis.

\begin{lstlisting}[language=XML, caption={Excerpt from the MJCF model file (\texttt{arm.xml}).}, label={lst:mjcf}]
<mujoco model="arm_9muscles">
  <option timestep="0.002" iterations="50" solver="Newton"
          integrator="implicit" cone="elliptic"/>

  <default>
    <muscle ctrllimited="true" ctrlrange="0 1"
            force="-1" scale="200" lmin="0.5" lmax="1.6"
            vmax="10" fpmax="1.5" fvmax="1.4"/>
  </default>

  <worldbody>
    <body name="humerus" pos="0 0 -0.3">
      <joint name="shoulder_flex" type="hinge" axis="0 1 0" range="-0.5 2.5" damping="0.1"/>
      <joint name="shoulder_abd"  type="hinge" axis="1 0 0" range="-0.5 2.5" damping="0.1"/>
      <joint name="shoulder_rot"  type="hinge" axis="0 0 1" range="-1.5 1.5" damping="0.1"/>
      <geom type="capsule" fromto="0 0 0 0 0 -0.3" size="0.04" mass="1.98"/>

      <!-- Muscle origin sites (humerus) -->
      <site name="bic_origin_long" pos="0     0.01  0.05"/>
      <site name="bic_s_origin"    pos="0.01  0.01  0.03"/>
      <site name="bra_origin"      pos="0     0.00 -0.15"/>
      <site name="tri_l_origin"    pos="0    -0.01  0.02"/>
      <site name="delt_a_origin"   pos="0.05  0.00  0.05"/>
      <site name="delt_m_origin"   pos="0.00  0.05  0.05"/>
      <site name="delt_p_origin"   pos="-0.05 0.00  0.05"/>

      <body name="radius" pos="0 0 -0.3">
        <joint name="elbow_flex" type="hinge" axis="0 1 0" range="0 2.7" damping="0.1"/>
        <geom type="capsule" fromto="0 0 0 0 0 -0.27" size="0.03" mass="1.18"/>

        <!-- Via-point and insertion sites (forearm) -->
        <site name="bic_via_1"        pos="0     0.02 -0.02"/>
        <site name="bic_insertion"    pos="0     0.02 -0.08"/>
        <site name="bic_s_insertion"  pos="0     0.02 -0.08"/>
        <site name="bra_insertion"    pos="0     0.01 -0.05"/>
        <site name="tri_l_insertion"  pos="0    -0.01 -0.01"/>
        <site name="tri_la_origin"    pos="0.01 -0.01 -0.10"/>
        <site name="tri_la_insertion" pos="0    -0.01 -0.01"/>
        <site name="tri_m_origin"     pos="-0.01 -0.01 -0.10"/>
        <site name="tri_m_insertion"  pos="0    -0.01 -0.01"/>
        <site name="delt_a_insertion" pos="0.03  0.00 -0.15"/>
        <site name="delt_m_insertion" pos="0.00  0.03 -0.15"/>
        <site name="delt_p_insertion" pos="-0.03 0.00 -0.15"/>

        <body name="hand" pos="0 0 -0.27">
          <geom type="sphere" size="0.04" mass="0.45"/>
          <site name="end_effector" pos="0 0 0"/>
        </body>
      </body>
    </body>
  </worldbody>

  <tendon>
    <spatial name="biceps_long"  width="0.003">
      <site site="bic_origin_long"/>
      <site site="bic_via_1"/>
      <site site="bic_insertion"/>
    </spatial>
    <spatial name="biceps_short" width="0.003">
      <site site="bic_s_origin"/>
      <site site="bic_s_insertion"/>
    </spatial>
    <spatial name="brachialis"   width="0.003">
      <site site="bra_origin"/>
      <site site="bra_insertion"/>
    </spatial>
    <spatial name="triceps_long" width="0.003">
      <site site="tri_l_origin"/>
      <site site="tri_l_insertion"/>
    </spatial>
    <spatial name="triceps_lat"  width="0.003">
      <site site="tri_la_origin"/>
      <site site="tri_la_insertion"/>
    </spatial>
    <spatial name="triceps_med"  width="0.003">
      <site site="tri_m_origin"/>
      <site site="tri_m_insertion"/>
    </spatial>
    <spatial name="deltoid_ant"  width="0.003">
      <site site="delt_a_origin"/>
      <site site="delt_a_insertion"/>
    </spatial>
    <spatial name="deltoid_med"  width="0.003">
      <site site="delt_m_origin"/>
      <site site="delt_m_insertion"/>
    </spatial>
    <spatial name="deltoid_post" width="0.003">
      <site site="delt_p_origin"/>
      <site site="delt_p_insertion"/>
    </spatial>
  </tendon>

  <actuator>
    <muscle name="biceps_long"  tendon="biceps_long"  force="624"  lengthrange="0.01 0.50"/>
    <muscle name="biceps_short" tendon="biceps_short" force="435"  lengthrange="0.01 0.50"/>
    <muscle name="brachialis"   tendon="brachialis"   force="987"  lengthrange="0.01 0.50"/>
    <muscle name="triceps_long" tendon="triceps_long" force="798"  lengthrange="0.01 0.50"/>
    <muscle name="triceps_lat"  tendon="triceps_lat"  force="624"  lengthrange="0.01 0.50"/>
    <muscle name="triceps_med"  tendon="triceps_med"  force="624"  lengthrange="0.01 0.50"/>
    <muscle name="deltoid_ant"  tendon="deltoid_ant"  force="1142" lengthrange="0.01 0.50"/>
    <muscle name="deltoid_med"  tendon="deltoid_med"  force="1142" lengthrange="0.01 0.50"/>
    <muscle name="deltoid_post" tendon="deltoid_post" force="259"  lengthrange="0.01 0.50"/>
  </actuator>

  <sensor>
    <actuatorfrc name="f_bl"  actuator="biceps_long"/>
    <actuatorfrc name="f_bs"  actuator="biceps_short"/>
    <actuatorfrc name="f_bra" actuator="brachialis"/>
    <actuatorfrc name="f_tl"  actuator="triceps_long"/>
    <actuatorfrc name="f_tla" actuator="triceps_lat"/>
    <actuatorfrc name="f_tm"  actuator="triceps_med"/>
    <actuatorfrc name="f_da"  actuator="deltoid_ant"/>
    <actuatorfrc name="f_dm"  actuator="deltoid_med"/>
    <actuatorfrc name="f_dp"  actuator="deltoid_post"/>
    <jointpos name="shf_q" joint="shoulder_flex"/>
    <jointvel name="shf_v" joint="shoulder_flex"/>
  </sensor>
</mujoco>
\end{lstlisting}

% -----------------------------------------------------------
\chapter[Gymnasium Environment (Python)]{Gymnasium Environment\\(Python)}
\label{app:env}

The Gymnasium environment \texttt{ArmReachEnv}, which integrates domain
randomisation and computes the reward of \cref{sec:reward}.

\begin{lstlisting}[language=Python, caption={Excerpt from \texttt{rl/environment.py}, as used for every reported result.}, label={lst:env}]
class ArmReachEnv(gym.Env):
    """Nine-muscle upper limb, 3D reaching task."""

    def __init__(self, domain_rand: bool = True, max_steps: int = None,
                 seed: int = None, reward_weights: dict = None, ...):
        self.model = mujoco.MjModel.from_xml_path(config.ARM_XML)
        self.data  = mujoco.MjData(self.model)

        # Nominal values kept so domain randomisation rescales from nominal
        # rather than compounding on the current value (see sec:defects).
        self._nominal_body_mass = self.model.body_mass.copy()
        self._nominal_gainprm   = self.model.actuator_gainprm.copy()
        self._nominal_damping   = self.model.dof_damping.copy()

        # Per-muscle peak isometric force, used to normalise the effort term.
        self._fmax = self._nominal_gainprm[:, 2].copy()

        self.action_space      = spaces.Box(0.0, 1.0, shape=(9,),  dtype=np.float32)
        self.observation_space = spaces.Box(-np.inf, np.inf, shape=(38,),
                                            dtype=np.float32)

        # Reward weights come from config.REWARD_WEIGHTS -- the single source
        # of truth. They are hand-set, NOT searched. An override dict is used
        # only by the reward-ablation study, to zero individual terms.
        w = dict(config.REWARD_WEIGHTS)
        if reward_weights:
            w.update(reward_weights)
        self.w1, self.w2, self.w3 = w["w1"], w["w2"], w["w3"]
        self.w4, self.w5          = w["w4"], w["w5"]
        self.w6                   = w.get("w6", 0.0)
        self.w7                   = w.get("w7", 0.0)

    def step(self, action):
        ...
        reward  = self.w4 - (self.w1 * (err**2 + 0.5 * err)
                             + self.w2 * e_eff
                             + self.w3 * float(np.mean(d_act**2)))
        reward += self.w7 * np.exp(-err / config.PRECISION_SCALE)
\end{lstlisting}

An earlier version of this appendix printed a hand-written excerpt that had
drifted from the code: it named the class \texttt{ArmMusculoEnv}, set
\texttt{max\_steps=1000} where the environment uses 100, gave the reward weights
as $(1.0, 0.005, 0.01)$, and annotated them ``optimised by Optuna''. The weights
are $(1.0, 1.0, 0.5)$ and were never searched --- \cref{sec:stats_protocol}
records that claim as withdrawn. The listing above is taken from the source that
produced the results.

% -----------------------------------------------------------
\chapter{SAC Training Script}
\label{app:train}

Training script using Stable-Baselines3 for \SAC{}, with
observation normalisation, periodic evaluation, and saving of
the best-performing models.

\begin{lstlisting}[language=Python, caption={Excerpt from \texttt{rl/train.py}.}, label={lst:train}]
import argparse, os, numpy as np
from stable_baselines3 import SAC
from stable_baselines3.common.vec_env import (
    SubprocVecEnv, VecNormalize)
from stable_baselines3.common.callbacks import EvalCallback
from arm_env import ArmMusculoEnv

def make_env(seed, domain_rand=True):
    def _init():
        env = ArmMusculoEnv(domain_rand=domain_rand, seed=seed)
        env.reset(seed=seed)
        return env
    return _init

def main(args):
    os.makedirs(args.logdir, exist_ok=True)

    # Environnements vectorises
    train_envs = SubprocVecEnv(
        [make_env(args.seed + i) for i in range(args.n_envs)])
    train_envs = VecNormalize(
        train_envs, norm_obs=True, norm_reward=True, clip_obs=10.0)

    eval_env = SubprocVecEnv([make_env(args.seed + 999)])
    eval_env = VecNormalize(
        eval_env, norm_obs=True, norm_reward=False,
        training=False)

    # SAC defaults as registered in rl/algorithms.py. The tuned run overrides
    # gamma (0.957), learning_rate (4.40e-4), batch_size (128) and tau (0.0188)
    # via the `hyperparams` passthrough; see tab:hparams.
    #
    # gradient_steps is set by the caller to 2 * config.N_ENVS = 8, giving a
    # measured replay ratio of 1.867. Leaving it at 1 is the defect of
    # sec:replaybug and must not be copied from here.
    model = SAC(
        "MlpPolicy", train_envs,
        learning_rate=3e-4,
        buffer_size=300_000,
        batch_size=512,
        learning_starts=4_000,
        tau=0.005,
        gamma=0.99,
        train_freq=1,
        gradient_steps=2 * config.N_ENVS,
        ent_coef="auto",
        target_entropy="auto",
        policy_kwargs=dict(net_arch=[256, 256]),
        seed=args.seed,
        verbose=1)
        # TensorBoard logging is deliberately not enabled (sec. 5.1).

    callback = EvalCallback(
        eval_env,
        best_model_save_path=args.logdir,
        log_path=args.logdir,
        eval_freq=5000,
        n_eval_episodes=20,
        deterministic=True)

    model.learn(total_timesteps=args.steps, callback=callback)
    model.save(os.path.join(args.logdir, "final_model"))
    train_envs.save(os.path.join(args.logdir, "vecnormalize.pkl"))

if __name__ == "__main__":
    parser = argparse.ArgumentParser()
    parser.add_argument("--seed", type=int, required=True)
    parser.add_argument("--steps", type=int, default=2_000_000)
    parser.add_argument("--n_envs", type=int, default=4)
    parser.add_argument("--logdir", type=str, default="runs/sac")
    args = parser.parse_args()
    main(args)
\end{lstlisting}

% -----------------------------------------------------------
\chapter[Hill Integration by Newton--Raphson]{Hill Integration\\by Newton--Raphson}
\label{app:newton}

Numerical solution of the muscle equilibrium equation~\eqref{eq:equilibrium}
by Newton--Raphson iterations, used at each simulation step to
determine the fibre length $\ell_M$ from the musculotendon length
$\ell_{MT}$ and the activation level $a$.

\begin{lstlisting}[language=Python, caption={Musculotendon equilibrium solver.}, label={lst:newton}]
import numpy as np

def solve_muscle_equilibrium(L_MT, a, L_M0, L_S0, alpha0,
                              F_max, epsilon=1e-6, max_iter=50):
    """
    Resout f_CE + f_PEE - f_SEE = 0 pour L_M.
    Retourne (L_M, L_S, F_muscle) par iterations de Newton-Raphson.
    """
    # Initialisation : hypothese L_S = L_S0 (tendon a longueur nominale)
    L_M = L_MT - L_S0
    L_M = max(0.5 * L_M0, min(1.6 * L_M0, L_M))  # bornes

    for it in range(max_iter):
        # Angle de pennation
        sin_alpha = (L_M0 * np.sin(alpha0)) / L_M
        cos_alpha = np.sqrt(max(1.0 - sin_alpha**2, 1e-6))

        # Longueur du tendon
        L_S = L_MT - L_M * cos_alpha

        # Force active (gaussienne centree sur L_M0)
        f_l = np.exp(-((L_M / L_M0 - 1.0) / 0.45) ** 2)
        f_CE = a * f_l * F_max

        # Force passive (exponentielle)
        if L_M > L_M0:
            f_PEE = F_max * (np.exp(5.0 * (L_M / L_M0 - 1.0)) - 1.0) \
                    / (np.exp(5.0) - 1.0)
        else:
            f_PEE = 0.0

        # Force tendineuse (quadratique en deformation)
        eps_s = (L_S - L_S0) / L_S0
        if eps_s > 0:
            f_SEE = F_max * 37.5 * eps_s ** 2
        else:
            f_SEE = 0.0

        # Equilibrium residual
        res = (f_CE + f_PEE) * cos_alpha - f_SEE
        if abs(res) < epsilon:
            break

        # Derivee numerique pour mise a jour
        dL = 1e-5 * L_M0
        L_M_p = L_M + dL
        # (recalcul rapide f_CE, f_PEE, f_SEE a L_M_p omis pour concision)
        # Puis : L_M -= res / (dRes / dL_M)
        L_M -= 0.5 * res / (F_max / L_M0)  # simplification pas adaptatif
        L_M = max(0.5 * L_M0, min(1.6 * L_M0, L_M))

    return L_M, L_S, (f_CE + f_PEE) * cos_alpha
\end{lstlisting}

% -----------------------------------------------------------
\chapter{Detailed Statistical Protocol}
\label{app:stats}

\textbf{This protocol was not applied to any result in this thesis, and no
$p$-value, effect size or confidence interval appears anywhere in the text.} It
is retained here because the implementation is complete and correct, and because
the seed count --- not the analysis --- is what prevents its use.

With the three seeds the available hardware permitted (\cref{sec:stats_protocol}),
the two-sided Wilcoxon signed-rank test has a minimum attainable $p$-value of
$0.25$. Running it would produce numbers that cannot reach significance under any
threshold while giving the appearance of inferential rigour. An earlier version of
this work reported $p$-values from this protocol over a purported ten seeds; those
analyses were never executed and are withdrawn (\cref{ch:conclusion}).

The code below therefore documents what \emph{should} be run once the replication
recommended in \cref{sec:future} is complete: Wilcoxon signed-rank test for paired
samples, Bonferroni correction for multiple comparisons, BCa bootstrap confidence
intervals, and rank-biserial effect size.

\begin{lstlisting}[language=Python, caption={Inter-algorithm statistical analysis.}, label={lst:stats}]
import numpy as np
from scipy import stats
from scipy.stats import wilcoxon

def rank_biserial(x, y):
    """Rank-biserial effect size for two paired samples."""
    diff = np.asarray(x) - np.asarray(y)
    pos = np.sum(diff > 0)
    neg = np.sum(diff < 0)
    return (pos - neg) / (pos + neg) if (pos + neg) > 0 else 0.0

def bootstrap_bca(data, stat_fn=np.mean, n_boot=10000,
                  alpha=0.05, rng=None):
    """BCa confidence interval (bias-corrected and accelerated)."""
    rng = np.random.default_rng(rng)
    data = np.asarray(data)
    n = len(data)
    boot_stats = np.array([
        stat_fn(rng.choice(data, size=n, replace=True))
        for _ in range(n_boot)])
    theta_hat = stat_fn(data)

    # Bias correction
    z0 = stats.norm.ppf(np.mean(boot_stats < theta_hat))
    # Acceleration via jackknife estimate
    jack = np.array([
        stat_fn(np.delete(data, i)) for i in range(n)])
    jack_mean = np.mean(jack)
    num = np.sum((jack_mean - jack) ** 3)
    den = 6.0 * (np.sum((jack_mean - jack) ** 2) ** 1.5 + 1e-12)
    a = num / den

    z_a = stats.norm.ppf(alpha / 2)
    z_1a = stats.norm.ppf(1 - alpha / 2)
    a1 = stats.norm.cdf(z0 + (z0 + z_a) / (1 - a * (z0 + z_a)))
    a2 = stats.norm.cdf(z0 + (z0 + z_1a) / (1 - a * (z0 + z_1a)))
    return (np.quantile(boot_stats, a1),
            np.quantile(boot_stats, a2))

def compare_algorithms(results):
    """results: dict {algo: array[10] de taux de succes}"""
    keys = list(results.keys())
    n_tests = len(keys) * (len(keys) - 1) // 2
    alpha_corrected = 0.05 / n_tests  # Bonferroni

    for i, a in enumerate(keys):
        for b in keys[i+1:]:
            stat, p = wilcoxon(results[a], results[b],
                               alternative='two-sided')
            rb = rank_biserial(results[a], results[b])
            ci_a = bootstrap_bca(results[a])
            print(f"{a} vs {b}: W={stat}, p={p:.4f}, "
                  f"rb={rb:.2f}, signif={p < alpha_corrected}")
\end{lstlisting}

% -----------------------------------------------------------
% CITATION FINALE
% -----------------------------------------------------------
\clearpage
\thispagestyle{empty}
\vspace*{\fill}
% An unattributed French epigraph stood here. It could not be
% traced to any source, so it was removed rather than left
% presented as a quotation.

